% ============================================================
% 纯答案册·课时12-2.5.2椭圆的几何性质 —— 工具/答案抽册器.py 抽册生成件（勿手改；第三档＝纯答案单独成册）
% 源件（只读）：C:/提示词/工作区/M3-第2章量产0913/成卷/导学件/课时12-2.5.2椭圆的几何性质/main.tex（md5 32a42af2fa2ad4ea3efc1e057435df4c）
%% 口径：题面不重印；逐题＝题号(印面号同正册)＋[答案]＋[详解]，值/详解照源件逐字
%       （钉值门硬断言：册值≡正册件内值≡值台账值tex，逐字全等）。册序＝正册装配序＝印面号 1..21 升序（同序同号）；键序断言基准＝题面库冻结 manifest canonical 键序。
%% 三档导出：整体 main-true／纯题 main-false（在产）｜本册＝纯答案册（本件）
%% 双档：ansbook-true.tex＝详解印本；ansbook-false.tex＝纯值速查本（详解不印）
% 编译：xelatex <壳> ×2，cwd＝本目录；sty＝qp-m3.sty 本地挂载（md5 7c3930362be8a0a2bdf21bbf8ac16573，锁校验）
% ============================================================
\documentclass[fontset=none]{ctexart}
\usepackage{qp-m3}
% —— 印面逐字符号路由补钉（承源件；− 走 TNR、▱ NSC 子块；禁改 sty 保 md5） ——
\xeCJKDeclareCharClass{Default}{"2212}%
\setCJKfamilyfont{nscblk}[Path=C:/提示词/工作区/字替对照-0909/variantF/fonts/,BoldFont=NSC-w600.ttf]{NSC-w500.ttf}%
\xeCJKDeclareSubCJKBlock{nscblk}{"25B1}%
\newcommand{\pxparallelogram}{{\CJKfamily{nscblk}▱}}%
% —— 断行微调（承母版实测档，三零门承重；\emergencystretch 不另设） ——
\xeCJKsetup{CJKglue={\hskip 0.10em plus 0.08em minus 0.05em}}%
% —— 纯答案册全线括线（长详解可跨栏断，承练习件全线括线口径；灰底盒不可跨栏断） ——
\ansblockgrayfalse
% —— 双档开关（false 档＝纯值速查：答案印、详解不印；件面开关，不动 sty 本体） ——
\newif\ifansbookdetail
\ifdefined\ansbookdetail
  \ifnum\ansbookdetail=0 \ansbookdetailfalse\else\ansbookdetailtrue\fi
\else
  \ansbookdetailtrue
\fi
% —— 页脚件名（承源件章名；件名＝<件型>答案册） ——
\renewcommand{\qpzhangming}{第二章\quad 平面解析几何}%
\renewcommand{\qpjianming}{导学件答案册}%

\begin{document}
\zhangtitle{第二章\quad 平面解析几何}
\jietitle{2.5.2 椭圆的几何性质}
\par\glueguard{1}\addvspace{2pt}\noindent\biaoqian{【纯答案册】}{\fangsong 题面不重印\quad 印面号与正册同号同序}\par\addvspace{2pt}

\begin{multicols}{2}
\emergencystretch=1em

\begin{ansblock}[2章-练-课时12-简1]
% ans:2章-练-课时12-简1
\ansitem{1}{x²/12＋y²/4＝1}
\ifansbookdetail
\ansnote{详解}{\(2a=4\sqrt{3}\)，即\(a=2\sqrt{3}\)，\(a^{2}=12\)；\(e=c/a=\sqrt{6}/3\)，即\(c^{2}=(6/9)\times 12=8\)，\(c=2\sqrt{2}\)；\(b^{2}=a^{2}-c^{2}=12-8=4\)．题设已限定焦点在\(x\)轴（方程形如\(x^{2}/a^{2}+y^{2}/b^{2}\)，\(a>b>0\)），故椭圆\(C\)的方程为\(x^{2}/12+y^{2}/4=1\)．}
\fi
\end{ansblock}

\begin{ansblock}[2章-练-课时12-简2]
% ans:2章-练-课时12-简2
\ansitem{2}{x²/3＋y²＝1（答案不唯一）}
\ifansbookdetail
\ansnote{详解}{设\(x^{2}/a^{2}+y^{2}/b^{2}=1\)（\(a>b>0\)），\(e^{2}=1-b^{2}/a^{2}=6/9=2/3\)，即\(a^{2}=3b^{2}\)．取\(b^{2}=1\)、\(a^{2}=3\)即得\(x^{2}/3+y^{2}=1\)；任取\(b^{2}=t\)（\(t>0\)）、\(a^{2}=3t\)均合题．}
\fi
\end{ansblock}

\begin{ansblock}[2章-练-课时12-简3]
% ans:2章-练-课时12-简3
\ansitem{3}{(−√2,√2)}
\ifansbookdetail
\ansnote{详解}{点在椭圆内部即代入左端小于1：\(m^{2}/4+1/2<1\)，即\(m^{2}/4<1/2\)，\(m^{2}<2\)，故\(-\sqrt{2}<m<\sqrt{2}\)．}
\fi
\end{ansblock}

\begin{ansblock}[2章-练-课时12-简4]
% ans:2章-练-课时12-简4
\ansitem{4}{A}
\ifansbookdetail
\ansnote{详解}{\(a^{2}=9\)、\(a=3\)，右焦点\(F_2(c,0)\)．椭圆上点到焦点距离的最小值在长轴（同侧）顶点取得：\(|PF_2|_{\min}=a-c\)．由\(3-c=3-2\sqrt{2}\)得\(c=2\sqrt{2}\)，\(b^{2}=a^{2}-c^{2}=9-8=1\)，\(b=1\)（\(b>0\)）．故选：A．}
\fi
\end{ansblock}

\begin{ansblock}[2章-练-课时12-简5]
% ans:2章-练-课时12-简5
\ansitem{5}{C}
\ifansbookdetail
\ansnote{详解}{由蒙日圆结论（本课时讲解位2.5.2.20）：\(x^{2}/a^{2}+y^{2}/b^{2}=1\)的两条互相垂直切线的交点轨迹为\(x^{2}+y^{2}=a^{2}+b^{2}\)．题设圆即该椭圆的蒙日圆，故\(6+b^{2}=8\)，\(b^{2}=2\)（满足\(0<b<\sqrt{6}\)）．故选：C．}
\fi
\end{ansblock}

\begin{ansblock}[2章-练-课时12-简6]
% ans:2章-练-课时12-简6
\ansitem{6}{(1)无数个，如 x²/3＋y²/8＝1、x²/5＋y²/10＝1（任写两个满足 a²−b²＝5 且焦点在 y 轴者均可）；(2)一个，x²/20＋y²/25＝1}
\ifansbookdetail
\ansnote{详解}{\(C\)化标准式\(x^{2}/4+y^{2}/9=1\)，焦点在\(y\)轴，\(c^{2}=9-4=5\)．(1)同焦点即\(a^{2}-b^{2}=5\)且焦点在\(y\)轴，即\(x^{2}/m+y^{2}/(m+5)=1\)（\(m>0\)）对任意\(m>0\)均可，故有无数个；取\(m=3\)、\(m=5\)得\(x^{2}/3+y^{2}/8=1\)、\(x^{2}/5+y^{2}/10=1\)．(2)代入\((4,\sqrt{5})\)：\(16/m+5/(m+5)=1\)，即\(16(m+5)+5m=m(m+5)\)，即\(m^{2}-16m-80=0\)，\((m-20)(m+4)=0\)，\(m=20\)或\(m=-4\)（舍，须\(m>0\)）．故唯一，方程\(x^{2}/20+y^{2}/25=1\)．}
\fi
\end{ansblock}

\begin{ansblock}[2章-练-课时12-简7]
% ans:2章-练-课时12-简7
\ansitem{7}{(−1,1)}
\ifansbookdetail
\ansnote{详解}{设\(A(x_1,y_1)\)、\(B(x_2,y_2)\)，“不关于长轴（\(x\)轴）对称”即\(y_1\neq -y_2\)（且\(A\)、\(B\)为不同两点）．\(|AM|=|BM|\)即\((x_1-m)^{2}+y_1^{2}=(x_2-m)^{2}+y_2^{2}\)，即\((x_1-x_2)(x_1+x_2-2m)=y_2^{2}-y_1^{2}\)．椭圆上\(y^{2}=12-(3/4)x^{2}\)，故\(y_2^{2}-y_1^{2}=(3/4)(x_1^{2}-x_2^{2})=(3/4)(x_1-x_2)(x_1+x_2)\)．若\(x_1=x_2\)，则由上式\(y_1^{2}=y_2^{2}\)，结合“不关于长轴对称”得\(y_1=y_2\)，\(A\)、\(B\)重合，矛盾；故\(x_1\neq x_2\)，可约去\((x_1-x_2)\)：\(x_1+x_2-2m=(3/4)(x_1+x_2)\)，即\(2m=(1/4)(x_1+x_2)\)，\(m=(x_1+x_2)/8\)．又\(-4\leqslant x_i\leqslant 4\)且\(x_1\neq x_2\)，故\(-8<x_1+x_2<8\)（两端点\(x_1=x_2=\pm 4\)不能同时取，等号不可达），得\(-1<m<1\)．}
\fi
\end{ansblock}

\begin{ansblock}[2章-练-课时12-简8]
% ans:2章-练-课时12-简8
\ansitem{8}{x²/4＋y²＝1}
\ifansbookdetail
\ansnote{详解}{\(P_3\)与\(P_4\)关于\(y\)轴对称，而椭圆关于\(y\)轴对称，故\(P_3\)、\(P_4\)同“在”或同“不在”；“恰有三点”须\(P_3\)、\(P_4\)都在其上，第三点为\(P_1\)、\(P_2\)之一．若\(P_1\)在\(C\)上，则与\(P_4\)同横坐标\(x=1\)而纵坐标不同（\(1\neq\sqrt{3}/2\)），同一\(x\)对应椭圆上两点纵坐标互为相反数，矛盾（亦可直接算：\(1/a^{2}+1/b^{2}=1\)与\(1/a^{2}+3/(4b^{2})=1\)相减得\(b=0\)，舍）；故第三点是\(P_2\)．由\(P_2(0,1)\)：\(1/b^{2}=1\)，即\(b^{2}=1\)；由\(P_4(1,\sqrt{3}/2)\)：\(1/a^{2}+3/4=1\)，即\(a^{2}=4\)（\(a>b\)成立）．所以\(C\)：\(x^{2}/4+y^{2}=1\)．回代检验：\(P_1(1,1)\)使\(1/4+1=5/4\neq 1\)，恰不在其上成立．}
\fi
\end{ansblock}

\begin{ansblock}[2章-练-课时12-简9]
% ans:2章-练-课时12-简9
\ansitem{9}{x＋y−1＝0}
\ifansbookdetail
\ansnote{详解}{\(a^{2}=4\)、\(b^{2}=3\)，即\(c^{2}=1\)，右焦点\(F(1,0)\)．设\(l\)的斜率为\(k\)（\(k\)必存在且\(k\neq 0\)：\(P\)在第二象限、\(F\)在\(x\)轴上），\(l\)：\(y=k(x-1)\)．设\(Q(x_2,y_2)\)，则\(Q'(x_2,-y_2)\)，且\(P\)、\(Q\)、\(F\)共线，故直线\(PQ\)即\(l\)，方向向量\((1,k)\)；向量\(FQ'=(x_2-1,-y_2)=(x_2-1)(1,-k)\)（用\(y_2=k(x_2-1)\)），即\(FQ'\)的方向为\((1,-k)\)（\(x_2\neq 1\)，否则\(Q\)在\(x\)轴上、\(Q\)为顶点使\(P\)不在第二象限，舍）．\(PQ\perp FQ'\)即\((1,k)\cdot(1,-k)=0\)，即\(1-k^{2}=0\)，\(k=\pm 1\)．由\(P\)在第二象限定号：\(k=1\)时\(l\)：\(y=x-1\)与\(x^{2}/4+y^{2}/3=1\)联立得\(7x^{2}-8x-8=0\)，两根之积\(-8/7<0\)，其负根对应\(y=x-1<0\)，点在第三象限，不合；\(k=-1\)时\(l\)：\(y=-(x-1)\)同得\(7x^{2}-8x-8=0\)，负根\(x=(4-6\sqrt{2})/7\)对应\(y=1-x>0\)，点在第二象限成立．故\(k=-1\)，\(l\)：\(y=-(x-1)\)，即\(x+y-1=0\)．}
\fi
\end{ansblock}

\begin{ansblock}[2章-练-课时12-简10]
% ans:2章-练-课时12-简10
\ansitem{10}{C}
\ifansbookdetail
\ansnote{详解}{同一中心天体（太阳）即\(a^{3}/T^{2}\)同值，故\(T_2^{2}=(a_2^{3}/a_1^{3})\times T_1^{2}=(60/1.5)^{3}\times 1^{2}=40^{3}=64000\)，\(T_2=\sqrt{64000}=80\sqrt{10}\approx 80\times 3.1=248\)（年）．故选：C．}
\fi
\end{ansblock}

\begin{ansblock}[2章-练-课时12-中1]
% ans:2章-练-课时12-中1
\ansitem{11}{[√5,3]}
\ifansbookdetail
\ansnote{详解}{\(|PA|+|PB|=6>|AB|=4\)，故\(P\)的轨迹是以\(A\)、\(B\)为焦点、\(2a=6\)的椭圆．以\(M\)为原点、\(AB\)所在直线为\(x\)轴建系，则\(A(-2,0)\)、\(B(2,0)\)，\(c=2\)、\(a=3\)、\(b^{2}=a^{2}-c^{2}=9-4=5\)，轨迹\(x^{2}/9+y^{2}/5=1\)．\(M\)即椭圆中心，\(|PM|\)为中心到椭圆上点的距离：\(|PM|^{2}=x^{2}+y^{2}=x^{2}+5(1-x^{2}/9)=5+(4/9)x^{2}\)，\(x^{2}\in[0,9]\)，故\(|PM|^{2}\in[5,9]\)，\(|PM|\in[\sqrt{5},3]\)（最小在短轴端点\(=b\)，最大在长轴端点\(=a\)）．}
\fi
\end{ansblock}

\begin{ansblock}[2章-练-课时12-中2]
% ans:2章-练-课时12-中2
\ansitem{12}{[4＋2√3, 8)}
\ifansbookdetail
\ansnote{详解}{\(a=2\)、\(b=\sqrt{3}\)、\(c=1\)，右焦点\(F(1,0)\)，记左焦点\(F_1(-1,0)\)．\(l\)过原点，\(A\)、\(B\)关于中心\(O\)对称，故\(O\)是\(AB\)的中点，又\(O\)是\(FF_1\)的中点，所以四边形\(AFBF_1\)是平行四边形，\(|BF|=|AF_1|\)．周长\(L=|AB|+|AF|+|BF|=|AB|+(|AF|+|AF_1|)=|AB|+2a=|AB|+4\)．求\(|AB|\)：\(l\)的斜率不为0（斜率为0时\(A\)、\(B\)为长轴端点，与\(F\)共线，三角形退化），设\(l\)：\(x=my\)，代入得\(y^{2}=12/(3m^{2}+4)\)，\(|AB|=2\sqrt{x^{2}+y^{2}}=2\sqrt{(m^{2}+1)y^{2}}=2\sqrt{12(m^{2}+1)/(3m^{2}+4)}=4\sqrt{1-1/(3m^{2}+4)}\)．\(m^{2}\geqslant 0\)，即\(3m^{2}+4\geqslant 4\)，\(1-1/(3m^{2}+4)\in[3/4,1)\)，故\(|AB|\in[2\sqrt{3},4)\)．所以\(L=|AB|+4\in[4+2\sqrt{3},8)\)（上端\(|AB|\to 2a=4\)对应直线趋于长轴、三角形退化，取开）．}
\fi
\end{ansblock}

\begin{ansblock}[2章-练-课时12-中3]
% ans:2章-练-课时12-中3
\ansitem{13}{最大 12(2＋√2)/5，最小 12(2−√2)/5}
\ifansbookdetail
\ansnote{详解}{设\(P(4\cos\theta,3\sin\theta)\)，则\(d=|3\cdot 4\cos\theta-4\cdot 3\sin\theta-24|/\sqrt{3^{2}+4^{2}}=|12\cos\theta-12\sin\theta-24|/5=|12\sqrt{2}\cos(\theta+\pi/4)-24|/5\)．因\(12\sqrt{2}<24\)，\(\cos(\theta+\pi/4)\in[-1,1]\)时\(12\sqrt{2}\cos(\theta+\pi/4)-24\)恒为负，绝对值直接翻号：\(d=(24-12\sqrt{2}\cos(\theta+\pi/4))/5\)．\(\cos(\theta+\pi/4)=-1\)时\(d_{\max}=(24+12\sqrt{2})/5=12(2+\sqrt{2})/5\)；\(\cos(\theta+\pi/4)=1\)时\(d_{\min}=(24-12\sqrt{2})/5=12(2-\sqrt{2})/5\)．（另法：设与\(3x-4y=24\)平行的切线\(3x-4y=t\)，由判别式为0得\(t=\pm 12\sqrt{2}\)，两切线到已知直线的距离\(|24\mp 12\sqrt{2}|/5\)即最小、最大值，结果一致．）}
\fi
\end{ansblock}

\begin{ansblock}[2章-练-课时12-中4]
% ans:2章-练-课时12-中4
\ansitem{14}{(1)2√3 km²；(2)24/7 km}
\ifansbookdetail
\ansnote{详解}{设平行四边形为\(ADBC\)（\(AB\)、\(CD\)为两条对角线），周长8 km即相邻两边之和\(|CA|+|CB|=4\) km，故另两顶点\(C\)（及\(D\)）在以\(A\)、\(B\)为焦点的椭圆上：\(2a=4\)、\(2c=|AB|=2\)，\(a=2\)、\(c=1\)、\(b^{2}=3\)．以\(AB\)所在直线为\(x\)轴、\(AB\)中点为原点建系，\(A(-1,0)\)、\(B(1,0)\)，椭圆\(x^{2}/4+y^{2}/3=1\)（\(y\neq 0\)，退化时不成四边形）．(1)平行四边形\(ADBC\)的面积\(=2S_{\triangle ABC}=2\times(1/2\times|AB|\times|y_C|)=2|y_C|\)，而\(C\)在椭圆上即\(|y_C|\leqslant b=\sqrt{3}\)，当\(C\)为短轴端点时取等，故最大面积\(=2\sqrt{3}\) km\(^{2}\)．(2)水沟过\(A(-1,0)\)且与\(AB\)（\(x\)轴）成45°角，取斜率1，方程\(y=x+1\)（取斜率\(-1\)时由对称性弦长相同）．“被农艺园围住的部分”即直线被椭圆截得的弦（口径已由亲算⑥注记2钉死，非“最大面积菱形”口径）：代入\(x^{2}/4+(x+1)^{2}/3=1\)，即\(3x^{2}+4(x^{2}+2x+1)=12\)，即\(7x^{2}+8x-8=0\)，\(x_1+x_2=-8/7\)，\(x_1x_2=-8/7\)，\(|x_1-x_2|=\sqrt{(8/7)^{2}+4\times 8/7}=\sqrt{288/49}=(12\sqrt{2})/7\)，弦长\(=\sqrt{1+k^{2}}\,|x_1-x_2|=\sqrt{2}\times(12\sqrt{2})/7=24/7\) km．}
\fi
\end{ansblock}

\begin{ansblock}[2章-练-课时12-难1]
% ans:2章-练-课时12-难1
\ansitem{15}{6}
\ifansbookdetail
\ansnote{详解}{\(A(0,b)\)、\(C(0,-b)\)，设\(B(x_1,y_1)\)、\(D(x_2,y_2)\)．\(\angle BAD=\angle BCD=90^{\circ}\)即\(A\)、\(C\)都在以\(BD\)为直径的圆上，即\(A\)、\(B\)、\(C\)、\(D\)四点共圆，圆心为\(BD\)的中点\(((x_1+x_2)/2,(y_1+y_2)/2)\)．又该圆过\(A\)、\(C\)，圆心在弦\(AC\)的垂直平分线（即\(x\)轴）上，故\((y_1+y_2)/2=0\)，即\(y_2=-y_1\)．面积条件：\(\triangle ABC\)与\(\triangle ADC\)共底\(AC\)（在\(y\)轴上），故\(S_{\triangle ABC}/S_{\triangle ADC}=|x_1|/|x_2|=2\)，即\(|x_1|=2|x_2|\)；因\(B\)、\(D\)分居\(y\)轴两侧，取\(x_1=-2x_2\)．于是圆心为\((x_1/4,0)\)，半径平方\(R^{2}=(x_1-x_1/4)^{2}+y_1^{2}=(9/16)x_1^{2}+y_1^{2}\)．把\(A(0,b)\)代入圆方程：\((x_1/4)^{2}+b^{2}=(9/16)x_1^{2}+y_1^{2}\)，即\(b^{2}=(1/2)x_1^{2}+y_1^{2}\)……①又\(B\)在椭圆上：\(x_1^{2}/18+y_1^{2}/b^{2}=1\)，配合①消\(x_1\)（\(x_1^{2}=2(b^{2}-y_1^{2})\)）：\((b^{2}-y_1^{2})/9+y_1^{2}/b^{2}=1\)，乘\(9b^{2}\)得\(b^{4}-b^{2}y_1^{2}+9y_1^{2}-9b^{2}=0\)，即\((b^{2}-9)(b^{2}-y_1^{2})=0\)．\(B\)异于上、下顶点，故\(y_1^{2}<b^{2}\)，\(b^{2}-y_1^{2}\neq 0\)，只能\(b^{2}=9\)，\(b=3\)，短轴长\(2b=6\)．（存在性核验：\(b^{2}=9\)时①化为\(x_1^{2}=18-2y_1^{2}\)，与椭圆\(x^{2}/18+y^{2}/9=1\)同解，构型可行．取\(y_1=0\)：\(B(3\sqrt{2},0)\)、\(D(-3\sqrt{2}/2,0)\)、\(A(0,3)\)、\(C(0,-3)\)，\(AB\cdot AD=(3\sqrt{2},-3)\cdot(-3\sqrt{2}/2,-3)=-9+9=0\)成立，\(\angle BAD=90^{\circ}\)，同理\(\angle BCD=90^{\circ}\)；\(S_{\triangle ABC}=1/2\times 6\times 3\sqrt{2}=9\sqrt{2}\)、\(S_{\triangle ADC}=1/2\times 6\times(3\sqrt{2}/2)=9\sqrt{2}/2\)，比值2成立．）}
\fi
\end{ansblock}

\begin{ansblock}[2章-练-课时12-难2]
% ans:2章-练-课时12-难2
\ansitem{16}{AD}
\ifansbookdetail
\ansnote{详解}{\(e=\sqrt{2}/2\)，即\(c^{2}/a^{2}=1/2\)，\(a^{2}=2b^{2}\)，\(c^{2}=a^{2}-b^{2}=b^{2}\)，即\(a=\sqrt{2}b\)、\(c=b\)．A：蒙日圆半径平方\(r^{2}=a^{2}+b^{2}=3b^{2}\)，圆心为中心，故\(x^{2}+y^{2}=3b^{2}\)，正确．B：把\(P(b,a)\)代入\(l\)：\(b\cdot b+a\cdot a-a^{2}-b^{2}=0\)，故\(P(b,a)\)在\(l\)上；又\(b^{2}+a^{2}=a^{2}+b^{2}=r^{2}\)，故\(P\)在蒙日圆上，过\(P\)的两条切线互相垂直，取\(A\)、\(B\)为两切点时\(PA\cdot PB=0\)，“任意一点恒有\(PA\cdot PB>0\)”不成立，错误．C：\(|AF_1|+|AF_2|=2a\)，故\(|AH|-|AF_2|=|AH|+|AF_1|-2a\geqslant d(F_1,l)-2a\)．\(d(F_1,l)=|b\cdot(-c)+a\cdot 0-a^{2}-b^{2}|/\sqrt{a^{2}+b^{2}}=(bc+a^{2}+b^{2})/\sqrt{a^{2}+b^{2}}=(b^{2}+2b^{2}+b^{2})/(\sqrt{3}b)=4b/\sqrt{3}=(4\sqrt{3}/3)b\)．故\(|AH|-|AF_2|\)的最小值为\((4\sqrt{3}/3)b-2a=(4\sqrt{3}/3)b-2\sqrt{2}b\)，并非\((4\sqrt{3}/3)b\)，错误（\((4\sqrt{3}/3)b\)只是\(F_1\)到\(l\)的距离，漏减\(2a\)）．D：矩形四边均与\(C\)相切，则相邻两边为一对互相垂直的切线，其四个顶点都在蒙日圆上，即蒙日圆是矩形的外接圆，对角线为直径\(2r=2\sqrt{3}b\)．设两邻边为\(p\)、\(q\)，则\(p^{2}+q^{2}=(2r)^{2}=12b^{2}\)，面积\(pq\leqslant(p^{2}+q^{2})/2=6b^{2}\)，当\(p=q=\sqrt{6}b\)（正方形、对角线落在坐标轴上）时取等，正确．故选：AD．}
\fi
\end{ansblock}

\begin{ansblock}[2章-导-课时12-G1]
% ans:2章-导-课时12-G1
\ansitem{17}{A}
\ifansbookdetail
\ansnote{详解}{\(e=\sqrt{1-b^{2}/a^{2}}\)，四个椭圆\(a^{2}\)同为20，故\(b^{2}\)越小\(e\)越大、椭圆越扁．\(b^{2}/a^{2}\)依次为\(9/20<10/20<11/20<12/20\)，所以\(x^{2}/20+y^{2}/9=1\)的离心率最大，形状最扁．故选：A．}
\fi
\end{ansblock}

\begin{ansblock}[2章-导-课时12-G2]
% ans:2章-导-课时12-G2
\ansitem{18}{B}
\ifansbookdetail
\ansnote{详解}{\(2c=|AB|=4\)，\(c=2\)；\(2a=|PA|+|PB|\)，\(e=c/a\)最大即\(a\)最小，即\(|PA|+|PB|\)最小．求折线和最小：\(A\)关于\(l\)：\(x-y+4=0\)的对称点\(A'(x',y')\)，由公式\(x'=x_0-2\cdot(x_0-y_0+4)/2=-2-2=-4\)，\(y'=y_0+2\cdot(x_0-y_0+4)/2=0+2=2\)，即\(A'(-4,2)\)（与“\(AA'\perp l\)、\(AA'\)中点在\(l\)上”两条件验算一致）．于是\(|PA|+|PB|=|PA'|+|PB|\geqslant|A'B|=\sqrt{(-4-2)^{2}+2^{2}}=\sqrt{40}=2\sqrt{10}\)，当\(P\)为\(A'B\)与\(l\)的交点时取等．故\(2a\geqslant 2\sqrt{10}\)，\(a\geqslant\sqrt{10}\)，\(e=c/a\leqslant 2/\sqrt{10}=\sqrt{10}/5\)．故选：B．}
\fi
\end{ansblock}

\begin{ansblock}[2章-导-课时12-G3]
% ans:2章-导-课时12-G3
\ansitem{19}{B}
\ifansbookdetail
\ansnote{详解}{\(e=\sqrt{6}/3\)，即\(c^{2}/a^{2}=2/3\)，\(b^{2}/a^{2}=1/3\)，\(a=\sqrt{3}b\)，直线\(ax-by=0\)即\(y=(a/b)x=\sqrt{3}x\)．圆\(M\)化为\((x-m/2)^{2}+y^{2}=(m^{2}-1)/4\)，圆心\((m/2,0)\)、半径\(r=\frac{1}{2}\sqrt{m^{2}-1}\)（须\(m^{2}>1\)）．相切即圆心到直线\(\sqrt{3}x-y=0\)的距离等于半径：\(|\sqrt{3}\cdot m/2|/2=\frac{1}{2}\sqrt{m^{2}-1}\)，平方得\(3m^{2}/16=(m^{2}-1)/4\)，即\(3m^{2}=4m^{2}-4\)，\(m^{2}=4\)，\(m=\pm 2\)（满足\(m^{2}>1\)）．故选：B．}
\fi
\end{ansblock}

\begin{ansblock}[2章-导-课时12-G4]
% ans:2章-导-课时12-G4
\ansitem{20}{(±√(5−m),0)（0＜m＜5）；(0,±√(m−5))（m＞5）}
\ifansbookdetail
\ansnote{详解}{须\(m>0\)且\(m\neq 5\)．当\(0<m<5\)时焦点在\(x\)轴，\(a^{2}=5\)、\(b^{2}=m\)，\(c^{2}=5-m\)，焦点为\((\pm\sqrt{5-m},0)\)；当\(m>5\)时焦点在\(y\)轴，\(a^{2}=m\)、\(b^{2}=5\)，\(c^{2}=m-5\)，焦点为\((0,\pm\sqrt{m-5})\)．}
\fi
\end{ansblock}

\begin{ansblock}[2章-导-课时12-G5]
% ans:2章-导-课时12-G5
\ansitem{21}{(1)√5/3；(2)√3/2}
\ifansbookdetail
\ansnote{详解}{(1)\(a:b=3:2\)，取\(a=3k\)、\(b=2k\)，\(c^{2}=a^{2}-b^{2}=5k^{2}\)，\(e=c/a=\sqrt{5}k/3k=\sqrt{5}/3\)．(2)短轴端点\((0,\pm b)\)与焦点\((c,0)\)构成等边三角形：底边\(2b\)，腰\(\sqrt{b^{2}+c^{2}}=a\)，故\(a=2b\)；\(c^{2}=a^{2}-b^{2}=3b^{2}\)，\(e=c/a=\sqrt{3}b/2b=\sqrt{3}/2\)．}
\fi
\end{ansblock}

% ---- 尾块（\tailfill[30mm] 定高参——钉档 §三.3：无参在平衡栏呈「内容尾随框」，贴栏底须定高参；实测无参致末页平衡 Overfull\vbox 12.99pt，定高 30mm 三零） ----
\tailfill[30mm]

\end{multicols}
\end{document}
